% Vietnamese translation for OI Public Library
% Source-PDF: IOI中国国家候选队论文/2009/李骥扬/线段跳表——跳表的一个拓展.pdf
\documentclass[11pt]{article}
\usepackage{oipl-vn}

\newcommand{\Oh}{\mathcal{O}}
\newcommand{\Path}{\operatorname{Path}}
\newcommand{\RevPath}{\operatorname{Rev\mbox{-}Path}}
\newcommand{\DPi}{\mathrm{DP}_i}
\newcommand{\SEGi}{\mathrm{SEG}_i}
\newcommand{\maxlev}{\mathit{maxlev}}
\newcommand{\figdesc}[2]{%
  \begin{center}
  \fbox{\begin{minipage}{0.88\linewidth}
  \small\textbf{Hình #1.} #2
  \end{minipage}}
  \end{center}
}

\title{Danh sách nhảy đoạn: một mở rộng của danh sách nhảy}
\author{Lý Ký Dương\\Trường Trung học số 2 Thạch Gia Trang}
\date{Luận văn Trại đông 2009}

\begin{document}
\maketitle

\origtitle{Danh sách nhảy đoạn: một mở rộng của danh sách nhảy}
\sourcepath{IOI China candidate team papers / 2009 / Li Jiyang / Segment Skip Lists}

\begin{abstract}
Bài viết chủ yếu giới thiệu một mở rộng của tôi đối với danh sách nhảy
\textit{skip list} (còn gọi là danh sách nhảy, Skip Lists, SL, do William
Pugh phát minh năm 1987): \term{danh sách nhảy đoạn} (\textit{Segment Skip
Lists}, SSL), cùng các chứng minh hiệu quả liên quan. Bài viết cũng phân tích
ngắn gọn ưu thế của danh sách nhảy và các mở rộng của nó trong thi tin học.
\end{abstract}

\paragraph{Từ khóa.}
Danh sách nhảy; cây cân bằng; danh sách nhảy đoạn; cây đoạn.

\section*{Dẫn nhập}

Phạm vi thảo luận trong bài này chỉ giới hạn ở danh sách nhảy ngẫu nhiên dựa
trên kỳ vọng (Skip Lists, SL), không xét danh sách nhảy tất định kiểu 1-2-3
(1-2-3 Deterministic Skip List, 1-2-3 DSL).\footnote{Chú thích này tương ứng
với chú thích mở đầu trong bản gốc.}

Trong thi tin học, cấu trúc dữ liệu là một phần quan trọng. Các bài toán dạng
này thường khảo sát năng lực duy trì và xử lý dữ liệu. Chúng hay yêu cầu ta
duy trì một cấu trúc tuyến tính có thứ tự và thực hiện trên đó các thao tác
như chèn, truy vấn. Những cách cài đặt thường dùng gồm cây cân bằng, danh sách
liên kết chia khối, cây Fenwick, v.v.

Danh sách nhảy là một cấu trúc dữ liệu mới, hỗ trợ các thao tác từ điển như
chèn, xóa và tìm kiếm trên bảng có thứ tự. Danh sách nhảy dựa trên ngẫu nhiên
hóa và xác suất; độ phức tạp kỳ vọng của các thao tác chèn, xóa, tìm kiếm đều
là $\Oh(\log N)$. Đồng thời, do bản chất vẫn là một bảng tuyến tính, nó có
tính chất dễ viết và dễ mở rộng. Nó có thể thay thế cây cân bằng để cài đặt
các thao tác từ điển. Hơn nữa, vì cấu trúc đơn giản và không cần thao tác
xoay, danh sách nhảy hỗ trợ lưu các phần khác nhau trên các môi trường khác
nhau và thực hiện thao tác song song.\footnote{Nội dung liên quan xem
\textit{Concurrent Maintenance of Skip Lists} của William Pugh.}

Tuy nhiên, danh sách nhảy cơ bản chỉ hỗ trợ các thao tác từ điển như chèn, xóa
và tìm kiếm, chứ không hỗ trợ thao tác trên đoạn. Năng lực của danh sách nhảy
không chỉ dừng ở đó: chỉ cần phân tích sâu hơn cấu trúc của nó, ta có thể rút
ra từ danh sách nhảy một cấu trúc cây nhị phân. Từ đó thêm thông tin đoạn để
thu được cây đoạn. Bài viết này sẽ nâng cấp các thao tác chèn và xóa để chúng
có thể duy trì thông tin đoạn.

Danh sách nhảy sau khi mở rộng được gọi là danh sách nhảy đoạn. Nó có thể thay
thế hoàn toàn nhiều loại cây cân bằng, hỗ trợ cả thao tác từ điển lẫn thao tác
trên đoạn. Đồng thời nó vẫn giữ các đặc điểm vốn có: đơn giản, dễ dùng, và còn
hỗ trợ những thao tác đặc biệt như xử lý tuần tự. Danh sách nhảy đoạn có thể
trở thành một cấu trúc dữ liệu lý tưởng trong thi tin học.

\tableofcontents

\section{Giới thiệu danh sách nhảy}

\subsection{Từ danh sách liên kết đến danh sách nhảy}

Danh sách liên kết, với tư cách một cấu trúc dữ liệu thông dụng và linh hoạt,
thường được dùng để duy trì dữ liệu tuyến tính. Danh sách liên kết thực hiện
chèn và xóa rất thuận tiện, nhưng chi phí tìm một số trong danh sách có thể
rất lớn, với độ phức tạp thời gian $\Oh(N)$. Ngược lại, nếu duy trì một mảng
có thứ tự, ta có thể tìm kiếm nhị phân trong thời gian $\Oh(\log N)$, nhưng để
duy trì tính có thứ tự của mảng thì chèn và xóa lại có độ phức tạp lớn, cỡ
$\Oh(N)$. Nếu có thể duy trì một danh sách liên kết có thứ tự đồng thời hỗ trợ
tìm kiếm nhị phân, ta sẽ hoàn thành các thao tác chèn, tìm kiếm và xóa trên
một tập dữ liệu quy mô $n$ trong thời gian $\Oh(\log N)$.

Trước hết hãy xét truy vấn trên một danh sách liên kết tĩnh. Một ý tưởng đơn
giản là thêm cho một số nút trong danh sách liên kết một con trỏ phụ, gọi là
con trỏ tăng tốc, để khi thích hợp ta có thể nhảy qua một số nút và tăng tốc
tìm kiếm.

\figdesc{1.1}{Danh sách liên kết có con trỏ tăng tốc: ngoài các cạnh nối tiếp
thông thường còn có vài cạnh dài giúp bỏ qua nhiều nút.}

Nhưng nếu giới hạn mỗi nút chỉ nối ra một con trỏ tăng tốc, dùng cấu trúc
tương tự danh sách liên kết chia khối, độ phức tạp tốt nhất ta đạt được chỉ là
$\Oh(\sqrt n)$. Có thể thu được kết quả tốt hơn không? Ta nới lỏng giới hạn
``mỗi nút có một con trỏ tăng tốc''. Dựa trên phép tương tự với tìm kiếm nhị
phân trong mảng tĩnh, ta có thể thu được cấu trúc sau.

\figdesc{1.2}{Một ``danh sách nhảy tưởng tượng'': các con trỏ tăng tốc nối tới
các nút giữa của những đoạn có thể xuất hiện trong quá trình tìm kiếm nhị
phân.}

Mỗi con trỏ tăng tốc nối ra từ một nút trỏ tới trung điểm của đoạn mà nút đó
có thể làm đầu trái trong quá trình tìm kiếm nhị phân. Cấu trúc này có thể tìm
kiếm dữ liệu trong thời gian $\Oh(\log N)$; quá trình tìm kiếm chính là tìm
kiếm nhị phân, mỗi lần dùng con trỏ tăng tốc để tìm nút giữa của đoạn hiện tại
rồi so sánh để xác định đoạn chính xác hơn và tiếp tục chia đôi.

Tuy nhiên, khi đối mặt với dữ liệu động, cấu trúc này rõ ràng không thể được
duy trì với chi phí nhỏ trong quá trình chèn và xóa.

Ngẫu nhiên hóa thường được dùng để giảm độ phức tạp thời gian của thuật toán.
Thêm tư tưởng ngẫu nhiên hóa vào hướng trên, ta thu được danh sách nhảy.
Danh sách nhảy gồm $\mathit{Lev}$ tầng danh sách liên kết và các con trỏ đi
xuống. Khi chèn dữ liệu $v$, ta ngẫu nhiên sinh ra chiều cao $h$, rồi chèn các
nút có khóa $v$ vào các tầng từ $1$ đến $h$; với các nút từ tầng $2$ đến tầng
$h$, thêm con trỏ đi xuống trỏ tới nút tầng 1 có cùng khóa. $h$ nút này tạo
thành một khối (\textit{block}) có chiều cao $h$.

\figdesc{1.3}{Một danh sách nhảy gồm nhiều tầng danh sách liên kết; các bản
sao của cùng một khóa được nối theo chiều dọc thành một khối.}

Danh sách nhảy có cấu trúc như vậy có độ phức tạp kỳ vọng của tìm kiếm là
$\Oh(\log N)$; xem phần 3.

\subsection{Mô tả hình thức cấu trúc danh sách nhảy}

Danh sách nhảy gồm nhiều danh sách liên kết
$L_1,L_2,\ldots,L_{\maxlev}$ và các con trỏ đi xuống, thỏa mãn các điều kiện
sau.

\begin{enumerate}
  \item $L_i$ là một danh sách liên kết, tăng đơn điệu từ trái sang phải, chứa
  hai phần tử đặc biệt $-\infty$ (nút bắt đầu) và $+\infty$ (nút kết thúc).

  \item $L_1$ chứa mọi phần tử, $L_i$ với $i\ge 2$ chứa một phần các phần tử,
  và với mọi $x\in L_i$ đều có
  $x\in L_j$ với mọi $1\le j<i\le \maxlev$. Nói cách khác,
  \[
    L_{\maxlev}\subseteq L_{\maxlev-1}\subseteq\cdots\subseteq L_2\subseteq L_1.
  \]

  \item Với mỗi nút bất kỳ trong $L_i$ ($i\ge 2$), có một con trỏ đi xuống
  trỏ tới nút trong $L_{i-1}$ có cùng khóa.
\end{enumerate}

Một số thuật ngữ dùng trong phần sau:

\begin{description}[leftmargin=2.6em,style=nextline]
  \item[Cạnh ngang] Con trỏ trong các danh sách liên kết ở từng tầng.
  \item[Cạnh dọc] Con trỏ đi xuống.
  \item[Khối] Một thành phần liên thông nếu chỉ xét các cạnh dọc; số nút trong
  thành phần đó gọi là chiều cao của khối.
  \item[Nút trước, nút sau] Nút đứng ngay trước hoặc ngay sau một nút trong
  danh sách liên kết nơi nó thuộc về.
  \item[Nút đáy] Nút nằm ở tầng thứ 1.
  \item[Đường tìm kiếm $\Path(v)$] Dãy nút đi qua khi tìm khóa $v$ trong danh
  sách nhảy.
  \item[Đường tìm kiếm ngược $\RevPath(v)$] Dãy thu được bằng cách đảo ngược
  $\Path(v)$.
\end{description}

\subsection{Các thao tác từ điển: tìm kiếm, chèn và xóa}

\subsubsection*{Tìm kiếm}

Quá trình tìm kiếm chính là quá trình xây dựng $\Path(v)$. Con trỏ $p$ bắt đầu
từ nút đầu $-\infty$ ở tầng trên cùng $L_h$. Mỗi lần, nếu khóa của nút sau nút
hiện tại nhỏ hơn hoặc bằng giá trị cần tìm, ta dịch $p$ sang phải đến nút sau;
ngược lại, ta đi theo cạnh dọc xuống tầng dưới. Khi không thể di chuyển nữa
(đang ở nút tầng đáy và khóa của nút sau lớn hơn giá trị cần tìm), nút mà $p$
trỏ tới chính là giá trị lớn nhất trong dữ liệu không vượt quá giá trị tìm
kiếm.

\figdesc{1.4}{Một quá trình tìm kiếm: đường màu xanh là đường tìm kiếm, các
cạnh màu đỏ là những cạnh đã thử so sánh nhưng không đi qua.}

\subsubsection*{Chèn}

Trước hết ngẫu nhiên hóa chiều cao của nút cần chèn. Cách sinh như sau. Mỗi
nút được chèn ít nhất vào tầng cơ sở, nên chiều cao ban đầu ít nhất là $0$.
Trong vòng lặp, mỗi lần có xác suất $p$ làm chiều cao tăng thêm $1$, và xác
suất $1-p$ thoát vòng lặp, từ đó thu được chiều cao ngẫu nhiên cần dùng. Để
lập trình thuận tiện, có thể đặt chiều cao tối đa $\maxlev$ và cố định duy trì
$\maxlev$ danh sách liên kết. Nếu chiều cao sinh ra lớn hơn $\maxlev$, ta gán
nó bằng $\maxlev$.

Giá trị của xác suất $p$ thường được chọn là $1/4$; phân tích cụ thể để ở phần
3. Mã giả:

\begin{verbatim}
Function random_height;
    p := 1 / 4;
    h := 0;
    while random() <= p do h := h + 1;
    // random() tra ve mot gia tri trong [0, 1)
    if h > maxlev then h := maxlev;
    return h;
\end{verbatim}

Sau đó chèn nút vào các tầng $1,\ldots,\mathit{random\_height}$ của danh sách
nhảy.

\figdesc{1.5.1}{Trạng thái trước khi chèn: các cạnh đỏ chỉ vị trí trên từng
tầng nơi nút mới sẽ được chen vào.}

\figdesc{1.5.2}{Trạng thái sau khi chèn: nút mới xuất hiện trên các tầng đã
chọn ngẫu nhiên, đồng thời các cạnh dọc nối các bản sao của nó.}

Khi cài đặt, trước tiên có thể tìm giá trị cần chèn để nhận được $\Path(v)$,
rồi chèn nút mới sau các nút của những tầng
$1,\ldots,\mathit{random\_height}$ được thăm cuối cùng trong $\Path(v)$. Nói
cách khác, ở mỗi tầng trong các tầng đó, chèn nút mới sau nút được thăm cuối
cùng.

\subsubsection*{Xóa}

Xóa tương tự chèn. Trước tiên tìm giá trị cần xóa để nhận được $\Path(v)$, rồi
xóa khỏi danh sách nhảy các nút có khóa $v$.

\figdesc{1.6.1}{Trước tiên tìm các nút cần xóa.}

\figdesc{1.6.2}{Xóa các nút và cạnh màu đỏ, rồi nối lại các cạnh màu xanh để
khôi phục các danh sách liên kết từng tầng.}

Mã liên quan trong bản gốc được đặt ở phụ lục \texttt{SkipList.Pas}.

\section{Cấu trúc và thao tác cơ bản của danh sách nhảy đoạn}

\subsection{Trích xuất cây đoạn}

Mục này trích xuất cây đoạn từ cấu trúc danh sách nhảy, rồi từ đó thu được
danh sách nhảy đoạn để duy trì thông tin đoạn. Với một danh sách nhảy cho
trước:

\figdesc{2.1}{Một danh sách nhảy thông thường.}

Ta định nghĩa:

\begin{description}[leftmargin=2.6em,style=nextline]
  \item[Cạnh thật] Trong các cạnh không trỏ tới $+\infty$, gồm mọi cạnh dọc và
  các cạnh ngang trỏ tới nút cao nhất của một khối.
  \item[Cạnh ảo] Mọi cạnh trỏ tới $+\infty$, cùng các cạnh ngang trỏ tới các
  nút nằm dưới nút cao nhất của một khối.
\end{description}

Nếu ẩn mọi cạnh ảo trong danh sách nhảy, ta thu được:

\figdesc{2.2}{Cây tìm kiếm nhị phân ẩn trong danh sách nhảy sau khi chỉ giữ
các cạnh thật.}

Ta nhận ra rằng các nút cùng các cạnh thật tạo thành một cây tìm kiếm nhị phân
(BST).

Tiếp tục, với mọi $v$, nếu thêm nhãn $v$ lên các nút trong $\Path(v)$, ta thu
được cấu trúc sau:

\figdesc{2.3.1}{Tiền thân của danh sách nhảy đoạn: các điểm rời rạc được đánh
dấu dọc theo các đường tìm kiếm.}

Ở đây ta đã có thể thấy bóng dáng của cây đoạn.

Tiến thêm một bước, thay các điểm rời rạc được đánh dấu bằng các đoạn liên
tục. Đồng thời thêm tầng thứ 0 bên dưới tầng thứ 1, cũng chứa tất cả khóa, và
thêm cho mọi nút tầng thứ 1 con trỏ đi xuống tới nút tương ứng ở tầng thứ 0.
Ta thu được:

\figdesc{2.3.2}{Phiên bản cuối cùng của danh sách nhảy đoạn: một cây đoạn ẩn
trong danh sách nhảy, có thêm tầng 0 để biểu diễn các lá.}

Đến đây, ta đã thu được một cây đoạn ẩn trong danh sách nhảy. Tôi gọi nó là
\term{danh sách nhảy đoạn} (\textit{Segment Skip List}, SSL).

\subsection{Mô tả hình thức danh sách nhảy đoạn}

Danh sách nhảy đoạn gồm nhiều danh sách liên kết
$L_0,L_1,L_2,\ldots,L_{\maxlev}$ và các con trỏ đi xuống, thỏa mãn các điều
kiện sau.

\begin{enumerate}
  \item $L_i$ là một danh sách liên kết, tăng đơn điệu từ trái sang phải, chứa
  hai phần tử đặc biệt $-\infty$ (nút bắt đầu) và $+\infty$ (nút kết thúc).

  \item $L_0$ và $L_1$ chứa mọi phần tử, $L_i$ với $i\ge 2$ chứa một phần các
  phần tử, và với mọi $x\in L_i$ đều có
  $x\in L_j$ với mọi $0\le j<i\le \maxlev$. Tức là
  \[
    L_{\maxlev}\subseteq L_{\maxlev-1}\subseteq\cdots\subseteq L_2
    \subseteq L_1=L_0.
  \]

  \item Với mỗi nút bất kỳ trong $L_i$ ($i\ge 1$), có một con trỏ đi xuống
  trỏ tới nút trong $L_{i-1}$ có cùng khóa.

  \item Mỗi nút có khóa nhỏ hơn $+\infty$ tương ứng với một đoạn nửa kín nửa
  hở. Biên trái của đoạn là khóa của nút đó. Giả sử trong quá trình tìm nút
  này, cạnh dọc cuối cùng đi qua là $u\to v$; khi đó biên phải của đoạn là
  khóa của nút sau $u$. Giá trị này không cần lưu thêm, vì trong mỗi quá trình
  tìm kiếm có thể thu được động.
\end{enumerate}

Ngoài việc định nghĩa các nút tầng 0 là nút đáy, các thuật ngữ đã nêu ở phần
trước về cơ bản không đổi.

Định nghĩa như vậy nhất định cho ra một cây đoạn, trong đó nút gốc biểu diễn
đoạn $[-\infty,+\infty)$. Điều này rất dễ chứng minh; quá trình chứng minh
được lược bỏ ở đây.

\subsection{Duy trì hai loại thông tin đoạn}

Tiếp theo là tìm cách dùng cây đoạn ẩn trong danh sách nhảy để duy trì thông
tin đoạn. Trước hết chia thông tin đoạn thành hai loại:

\begin{description}[leftmargin=2.6em,style=nextline]
  \item[Thông tin DP ($\DPi$)] Thông tin của một đoạn được lưu khi thực hiện
  quy hoạch động trên cây, chẳng hạn số điểm có khóa nằm trong đoạn, giá trị
  lớn nhất, giá trị nhỏ nhất. Đặc trưng là mỗi nút đều lưu thông tin này.
  Thông tin có thể được tính từ thông tin của hai con của nút, thường trong
  thời gian $\Oh(1)$.

  \item[Thông tin đoạn ($\SEGi$)] Thông tin được lưu trong cây đoạn, chẳng hạn
  trạng thái tô màu của một đoạn. Đặc trưng là không phải mọi nút đều lưu
  thông tin này. Khi thực hiện thao tác đoạn, nếu đoạn thao tác chứa đoạn
  tương ứng với một nút không phải lá, có thể lưu thông tin tại nút không phải
  lá đó để giảm độ phức tạp. Thông tin cũng có thể được đẩy xuống hai nút con;
  ví dụ màu của một đoạn có thể được phân phối cho hai con.
\end{description}

Hai loại thông tin này có đặc điểm khác nhau, vì vậy phương pháp thao tác cũng
khác nhau. Sau đây giới thiệu cách danh sách nhảy đoạn duy trì hai loại thông
tin này khi chèn và xóa.

\subsubsection*{Chèn}

Mặc dù sau khi chèn nút, cấu trúc cây đoạn có thể thay đổi, nhưng như hình
dưới cho thấy, số nút có đoạn tương ứng bị thay đổi chỉ ở quy mô
$\Oh(\log N)$.

\figdesc{2.4.1}{Sự thay đổi các đoạn do nút SSL biểu diễn trong quá trình
chèn. Các nút viền xanh là nút mới thêm; các đoạn chữ đỏ là những đoạn bị
thay đổi. Ngoài một phần của $\Path(v)$ và $\Path(v-\varepsilon)$, đoạn tương
ứng với các nút khác không đổi.}

Vì vậy số nút cần cập nhật thông tin cũng chỉ có quy mô $\Oh(\log N)$.

Thứ tự thao tác khi chèn $v$:

\begin{enumerate}
  \item Dọc theo $\Path(v)$, đẩy thông tin $\SEGi$ xuống, và thông qua quá
  trình này thu được thông tin $\SEGi$ của nút $v$.
  \item Chèn khối có khóa $v$, đồng thời đặt giá trị $\SEGi$ của nút đáy bằng
  giá trị $\SEGi$ vừa tìm được.
  \item Theo thứ tự $\RevPath(v)$, cập nhật thông tin $\DPi$ của các nút trên
  đường.
  \item Theo thứ tự $\RevPath(v-\varepsilon)$, cập nhật thông tin $\DPi$ của
  các nút trên đường.
\end{enumerate}

Chỉ cần chứng minh hai kết luận sau là đủ để đảm bảo tính đúng đắn:

\begin{enumerate}
  \item Mọi nút thuộc $\Path(v)$ đều biểu diễn đoạn chứa $v$.
  \item Sau khi chèn $v$, chỉ các nút trong $\Path(v)$ và
  $\Path(v-\varepsilon)$ có phạm vi đoạn biểu diễn bị thay đổi.
\end{enumerate}

Chứng minh như sau. Với kết luận thứ nhất, mỗi lần đi tới một nút thì khóa của
nút đó nhất định nhỏ hơn hoặc bằng mục tiêu tìm kiếm. Do đó mọi điểm trong
$\Path(v)$ có khóa nhỏ hơn hoặc bằng $v$, tức biên trái đóng của đoạn do mỗi
điểm biểu diễn nhất định không vượt quá $v$. Với biên phải của đoạn, nó nhất
định xuất hiện tại một lần so sánh nào đó: sau khi phát hiện $v$ nhỏ hơn giá
trị này, ta không đi sang phải nữa mà đi xuống. Vì thế biên phải mở của đoạn
nhất định lớn hơn $v$.

Với kết luận thứ hai, trước hết khóa của từng điểm không đổi, nên biên trái
của mọi đoạn không đổi. Vì vậy chỉ cần xét biên phải. Theo định nghĩa, biên
phải là khóa của nút sau nút xuất phát của cạnh dọc cuối cùng trong quá trình
tìm kiếm. Với các nút nằm ngoài $\Path(v)$ và $\Path(v-\varepsilon)$, con trỏ
ngang nối ra từ nút xuất phát của cạnh dọc đó không thay đổi. Tổng hợp lại,
sau khi chèn $v$, chỉ các nút trong $\Path(v)$ và $\Path(v-\varepsilon)$ có
phạm vi đoạn biểu diễn thay đổi.

Những nút cần thay đổi thông tin đoạn chỉ là các nút có phạm vi đoạn bị đổi và
các nút có đoạn tương ứng chứa $v$, tức chỉ các điểm trên $\Path(v)$ và
$\Path(v-\varepsilon)$.

\subsubsection*{Xóa}

Xóa tương tự chèn. Đại thể quá trình như sau:

\begin{enumerate}
  \item Dọn các thông tin $\SEGi$ sẽ bị che phủ. Dọc theo $\Path(v)$, sau khi
  lần đầu gặp $\SEGi\ne \mathrm{NULL}$, đặt $\SEGi$ của mọi nút phía sau trên
  đường này thành rỗng. Sau đó lặp lại một lần dọc theo
  $\Path(v-\varepsilon)$.

  \item Nếu trong $\Path(v)$, thông tin $\SEGi$ của nút có khóa $v$ không hoàn
  toàn rỗng, thì $v$ là đầu trái của một đoạn tô màu nào đó và không thể xóa.
  Tiếp tục xét trong $\Path(v-\varepsilon)$ các nút có số tầng không lớn hơn
  chiều cao của khối có khóa $v$: nếu thông tin $\SEGi$ của chúng đều rỗng,
  thì $v$ là đầu phải của một đoạn nào đó và không thể xóa. Với những thông
  tin $\SEGi$ có thể gộp được, trước hết nên gộp ra giá trị trên $\Path(v)$ và
  $\Path(v-\varepsilon)$. Nếu không rơi vào các trường hợp trên thì tiếp tục.

  \item Xóa khối có khóa $v$.

  \item Dọc theo $\RevPath(v)$, cập nhật thông tin $\DPi$.
\end{enumerate}

Để rút gọn độ dài, ở đây tôi viết chung thao tác duy trì hai loại thông tin.
Nhưng trong ứng dụng thực tế thường chỉ dùng một loại thao tác. Phụ lục bản
gốc cung cấp riêng hai chương trình làm tham khảo cho hai loại thông tin:
\texttt{SEGiSSL.Pas} và \texttt{DPiSSL.Pas}.

\section{Ưu thế và phân tích hiệu quả của danh sách nhảy}

\subsection{Phân tích hiệu quả của danh sách nhảy}

\subsubsection*{Độ phức tạp thời gian kỳ vọng của tìm kiếm, chèn và xóa}

Trước hết giả sử $\maxlev\to+\infty$. Ta suy ra độ phức tạp kỳ vọng của tìm
kiếm, tức độ dài kỳ vọng của đường tìm kiếm.\footnote{Mục này đưa ra một suy
diễn ngắn gọn, trực quan nhưng chưa hoàn toàn nghiêm ngặt. Chứng minh chặt chẽ
xem \textit{Skip Lists: A Probabilistic Alternative to Balanced Trees} của
William Pugh.}

Gọi $f(N)$ là độ dài kỳ vọng của đường tìm kiếm trong danh sách nhảy có $N$
nút. Nó gồm ba phần:

\begin{enumerate}
  \item Độ dài kỳ vọng của đường tìm kiếm trong danh sách nhảy tạo bởi các
  tầng từ 1 đến $\maxlev$. Phần này tương đương với độ dài kỳ vọng của một
  danh sách nhảy có quy mô kỳ vọng $\Oh(pN)$.

  \item Một con trỏ đi từ tầng thứ hai xuống tầng thứ nhất. Trong danh sách
  nhảy đoạn, đây là con trỏ đi từ tầng thứ nhất xuống tầng thứ không.

  \item Độ dài đường đi sang phải ở tầng thứ nhất. Trong danh sách nhảy đoạn,
  đây là độ dài đi sang phải ở tầng thứ không. Mỗi lần có thể đi sang phải với
  xác suất $1-p$, ký hiệu là $q$.
\end{enumerate}

Do đó
\[
\begin{aligned}
  f(N)
    &= f(pN)+1+q+q^2+q^3+q^4+q^5+\cdots\\
    &= f(pN)+\frac{1}{p}.
\end{aligned}
\]

Giải hàm truy hồi
\[
  f(N)=f(pN)+\frac{1}{p}
\]
ta được
\[
  f(N)=\log_a N,
  \qquad
  a=\left(\frac{1}{p}\right)^p.
\]

Độ phức tạp thời gian của chèn và xóa gồm một phần liên quan tới độ dài đường
tìm kiếm và một phần liên quan tới chiều cao, tức
\[
  \Oh(\log_a N+\mathit{max\_height})=\Oh(\log_a N),
\]
trong đó $\mathit{max\_height}$ là chiều cao của danh sách nhảy. Khi $\maxlev$
là một hằng số, không tiến tới vô cùng, ta có
$\mathit{max\_height}=\maxlev$. Do bị giới hạn bởi $\maxlev$, độ phức tạp thời
gian trong tình huống này không dễ phân tích, nhưng trong thực tế không gây
suy biến, trừ khi đặt $\maxlev$ quá nhỏ. Nó làm việc lập trình đơn giản hơn,
vì có thể cố định số tầng của danh sách nhảy là $\maxlev$; cách này được
khuyến nghị. Cách chọn $\maxlev$ xem bên dưới.

\subsubsection*{Độ phức tạp không gian kỳ vọng}

Vẫn giả sử $\maxlev\to+\infty$.

\begin{enumerate}
  \item Kỳ vọng chiều cao của một nút đơn lẻ:
  \[
    h_x=1+p+p^2+p^3+p^4+\cdots=\frac{1}{1-p}.
  \]

  \item Độ phức tạp không gian kỳ vọng:
  \[
    \text{số nút}=N\times h_x=\frac{N}{1-p}.
  \]

  \item Phân tích chiều cao lớn nhất:
  \[
    \Pr(h_{\max}\le m)=\left(1-p^{m-1}\right)^N.
  \]

  \item Khi $\maxlev$ là một hằng số, không tiến tới vô cùng, kỳ vọng chiều
  cao của nút không vượt quá giá trị phân tích ở trên, và độ phức tạp không
  gian cũng không vượt quá giá trị thu được ở trên.
\end{enumerate}

\subsubsection*{Xác suất $p$}

Với
\[
  a=\left(\frac{1}{p}\right)^p,
\]
ta có
\[
  f(N)=\log_a N
      =\frac{\log_2 N}{\log_2\left((1/p)^p\right)}.
\]

Khi $p=1/e$, hệ số là nhỏ nhất:
\[
  f(N)
  =\frac{\log_2 N}{\log_2(e^{1/e})}
  \approx 1.88416939\log_2 N.
\]

Với $p=1/2$ và $p=1/4$, hệ số đều là
\[
  f(N)
  =\frac{\log_2 N}{\log_2(2^{1/2})}
  =\frac{\log_2 N}{\log_2(4^{1/4})}
  =2\log_2 N.
\]

$p=1/e$ có thời gian ngắn hơn $p=1/2$ và $p=1/4$ một chút, nhưng khó cài đặt
hơn, nên thông thường dùng $p=1/2$ hoặc $p=1/4$. Vì lý do không gian, $p=1/4$
được dùng phổ biến hơn.

Hệ số không gian là
\[
  b=\frac{1}{1-p}.
\]
Từ đó biết được khi $p=1/4$, lượng không gian dùng ít hơn $p=1/2$; xấp xỉ
bằng $2/3$ so với khi $p=1/2$. Vì vậy giá trị $p$ thường dùng nhất là $1/4$.

\subsubsection*{Giới hạn số tầng $\maxlev$}

Khi không có giới hạn tầng cao nhất, tức $\maxlev\to+\infty$, mới là danh sách
nhảy chuẩn. Nhưng trong ứng dụng thực tế, để cài đặt đơn giản, ta thường đặt
một hằng số như vậy. Sau khi đưa vào $\maxlev$, độ phức tạp thời gian không dễ
phân tích, nhưng hiệu quả thực nghiệm tốt. Xem phần phân tích chiều cao lớn
nhất: có thể chọn một $\maxlev$ phù hợp. Ví dụ với $p=1/4$, khi $N=10^6$ chọn
$\maxlev=16$, ta có
\[
  \Pr(h_{\max}\le \maxlev)
  =\left(1-p^{\maxlev-1}\right)^N
  \approx 0.999069111,
\]
là một lựa chọn không tệ.

\subsection{Phân tích hiệu quả của danh sách nhảy đoạn}

Các thao tác của danh sách nhảy đoạn về cơ bản giống danh sách nhảy, hiệu quả
cũng phụ thuộc vào độ dài kỳ vọng của đường tìm kiếm. Độ phức tạp thời gian
vẫn là $\Oh(\log N)$, chỉ có hằng số lớn hơn.

Về độ phức tạp không gian, do nó nhiều hơn danh sách nhảy thường một tầng, có
thể ghi thô là
\[
  \Oh\!\left(\frac{2-p}{1-p}N\right).
\]
Thoạt nhìn khá lớn, nhưng thực ra khi $p=1/4$ thì chỉ là
\[
  \Oh\!\left(\frac{7}{3}N\right),
\]
chỉ lớn hơn cây đoạn một chút. Do mỗi nút cần ghi thông tin đoạn, kích thước
bản thân nút cũng tăng. Nhưng độ phức tạp không gian vẫn là $\Oh(N)$.

\subsection{Ưu thế của danh sách nhảy}

\begin{itemize}
  \item Bản chất của danh sách nhảy là danh sách liên kết, thao tác đơn giản,
  không cần thao tác xoay. Danh sách nhảy ngắn hơn nhiều cây cân bằng thông
  thường. Quan trọng hơn, các thao tác của nó dễ suy nghĩ, không cần kỹ thuật
  viết đặc biệt để rút ngắn mã.

  \item Hỗ trợ thao tác song song. Danh sách nhảy không cần xoay, cấu trúc
  không sinh ra thay đổi lớn, có thể lưu trên các môi trường khác nhau và hỗ
  trợ thao tác song song. Xem \textit{Concurrent Maintenance of Skip Lists}
  của William Pugh.

  \item Chỉ cần thay đổi giá trị $p$ là có thể điều chỉnh hiệu quả thời gian
  và không gian.

  \item Danh sách nhảy nguyên thủy, chỉ hỗ trợ thao tác từ điển, không cần lưu
  bất kỳ thông tin phụ nào để duy trì cân bằng mà vẫn đạt độ phức tạp thời
  gian tốt. Không có lãng phí không gian.

  \item Có thể thực hiện các thao tác đặc biệt như thao tác tuần tự, tức xử lý
  các nút theo thứ tự khóa.
\end{itemize}

\subsection{So sánh danh sách nhảy đoạn với cây đoạn và cây cân bằng}

Cây đoạn và cây cân bằng là các cấu trúc dữ liệu thường dùng trong thi tin
học, rất nhiều bài toán có thể dùng chúng. Danh sách nhảy đoạn được giới thiệu
trong bài này dùng ngẫu nhiên hóa để thu được một cấu trúc tương đối cân bằng,
đồng thời có thể hỗ trợ mọi chức năng của cây đoạn và cây cân bằng. So với cây
đoạn, nó có đặc điểm là có thể xử lý thông tin động và hỗ trợ truy vấn trực
tuyến. So với cây cân bằng, nó có đặc điểm là dễ viết và mã ngắn gọn.

Tóm lại, danh sách nhảy và mở rộng của nó, danh sách nhảy đoạn, là các cấu
trúc dữ liệu xuất sắc. Chúng có thể trở thành một công cụ sắc bén cho các thí
sinh tham gia thi tin học.

\section*{Tài liệu tham khảo}

\begin{enumerate}
  \item William Pugh, \textit{Skip Lists: A Probabilistic Alternative to
  Balanced Trees}.
  \item William Pugh, \textit{Concurrent Maintenance of Skip Lists}.
  \item Wei Ran, luận văn Trại đông 2005, \textit{Làm hiệu quả thuật toán
  ``nhảy lên''! -- Bàn sơ về các thao tác liên quan và ứng dụng của danh sách
  nhảy}.
\end{enumerate}

\end{document}
